% !TeX program = xelatex
\begin{filecontents*}{audit-references.bib}
@book{liu2013,
  author = {刘海洋},
  title = {{LaTeX} 入门},
  address = {北京},
  publisher = {电子工业出版社},
  year = {2013}
}
@online{amsmathGuide,
  author = {{American Mathematical Society} and {The LaTeX Project}},
  title = {User's Guide for the {amsmath} Package},
  year = {2020},
  url = {https://tug.ctan.org/macros/latex/required/amsmath/amsldoc.pdf},
  urldate = {2026-11-04}
}
\end{filecontents*}
\documentclass[UTF8,11pt,a4paper,oneside,openany,linespread=1]{ctexbook}
\usepackage[margin=25mm]{geometry}
\usepackage{amsmath,amssymb,bm,booktabs,array,enumitem}
\usepackage{setspace}
\setstretch{1.28}
\setlength{\parskip}{3pt}
\setlength{\jot}{5pt}
\raggedbottom
\usepackage{gbt7714}
\usepackage[unicode,hidelinks]{hyperref}
\bibliographystyle{gbt7714-numerical}
\begin{document}
\setcounter{chapter}{7}
\chapter{数模论文实战（下）——文献与排错}

\section{中文文献与引用}
我用 gbt7714 排顺序编码制参考文献。示例沿用 TeX Live 2024 中的
\texttt{gbt7714-numerical}，对应 GB/T 7714—2015。

使用 gbt7714 3.1.1 时，将样式名改为 \texttt{gbt7714-numeric}。
GB/T 7714—2025 已于 2026 年 7 月实施，交稿前请核对接收方要求的版本。
\href{https://std.samr.gov.cn/gb/search/gbDetailed?id=4507EFE13D37CB6AE06397BE0A0A601F}{国家标准信息}

例子引用了刘海洋的中文入门书\cite{liu2013}和 amsmath 手册\cite{amsmathGuide}。
数据库中的引用键用来调用条目，不必自己手写文献编号。

filecontents 首次运行会建立演示数据库。若同名文件已经存在，
修改条目后要同步更新对应的 .bib 文件，或在新的演示文件夹中编译。
署名、标题和出版信息仍应从原始资料核对，排版工具不会替我查证。

\section{按四步编译}
这份例子使用传统 BibTeX，流程仍是四步：
\begin{enumerate}
  \item XeLaTeX 读取正文，写出引用和样式信息。
  \item BibTeX 读取辅助文件与数据库，生成文献内容。
  \item XeLaTeX 排出文献表，并将条目对应的编号写入辅助文件。
  \item XeLaTeX 再跑一次，使引用稳定。
\end{enumerate}
主文件若叫 main.tex，BibTeX 的作业名就是 main，不是数据库名。
gbt7714 与 biblatex-gb7714 不要混用；后者常搭配 Biber。

\section{三种常见报错，我怎样排查}
\subsection{少了括号，后面都成了嫌疑人}
看到 \verb|Missing } inserted| 时，我会查前面的花括号和分式。
少一个右花括号，后面几段就可能被吃进同一个参数。

看到 \verb|Runaway argument?| 时，我会先查上一条命令的参数是否完整。
再逐个核对相关环境的起止。

\subsection{命令拼错，也可能是宏包缺席}
Undefined control sequence 指向不认识的命令。
我先核对拼写、大小写和反斜杠，再查所需宏包或自定义命令。
拿一份最小例子验证，比一次塞进十个宏包可靠。

\subsection{文件明明在，为什么还说找不到}
File not found 时，我查编译工作目录、相对路径、扩展名和大小写，
也留意系统隐藏的双后缀。把 JPG 改成 PNG 后缀并不会转换格式。
日志是案发现场，最后那句“没生成 PDF”通常只是结案通知，线索还在前面。

\section{交稿前的核对表}
\begin{center}
\begin{tabular}{ll}
\toprule
对象 & 我会检查 \\
\midrule
首页 & 标题、页码、匿名信息与元数据 \\
引用 & 是否还有问号、跳转目标是否正确 \\
图表 & 单位、图题、数据来源和小数精度 \\
程序 & 是否能复现正文中的结果 \\
文献 & 署名、年份、来源和实际引用位置 \\
\bottomrule
\end{tabular}
\end{center}
我会实际打开 PDF 检查，而不只看源码里“应该没问题”。
Overfull 是内容越界的提醒，文件能生成也仍需修版面。
\bibliography{audit-references}
\end{document}
